% Part: computability
% Chapter: recursive-functions
% Section: bounded-minimization

\documentclass[../../../include/open-logic-section]{subfiles}

\begin{document}

\olfileid{cmp}{rec}{bmi}
\olsection{Minimisasi Mawa Wates}

\begin{explain}
Kadhangkala migunani yen kita netepake fungsi minangka wilangan
paling cilik kang nyukupi sipat utawa relasi~$P$. Yen $P$
bisa diputusake, kita bisa ngitung fungsi iki kanthi nyoba
kabeh wilangan kang mungkin, $0$, $1$, $2$, \dots, nganti
nemokake wilangan paling cilik kang nyukupi~$P$. Cathetan cakupan SOL6-F036: panggolekan tanpa wates iki rampung yen ana saksi; yen ora ana, tes kang bisa diputusake ora dhewe njamin panggolekan rampung.
Panggolekan tanpa wates kaya mangkene nggawa kita metu
saka golongan fungsi rekursif primitif. Nanging yen kita
mung nggoleki wilangan paling cilik kang \emph{kurang saka
wates kang diwenehake kanthi mandiri}, asil kasebut tetep
rekursif primitif yen relasi kang dipriksa uga rekursif primitif. Yen relasine mung bisa diputusake, wates winates njamin panggolekan rampung, nanging durung njamin asil rekursif primitif. Kanthi tembung liya, lan luwih umum,
upamana kita nduweni relasi rekursif primitif~$R(x,z)$.
Gatekna fungsi kang nggandhengake $x$ lan~$y$ karo
wilangan $z < y$ paling cilik kang nyukupi $R(x, z)$.
Fungsi iki uga bisa diitung kanthi mriksa apa $R(x, 0)$,
$R(x, 1)$, \dots, $R(x, y-1)$ bener. Nanging ngapa
fungsi iki rekursif primitif?
\end{explain}

\begin{prop}
Yen $R(\vec x, z)$ rekursif primitif, fungsi
$m_R(\vec{x}, y)$ kang ngasilake~$z$ paling cilik
kang kurang saka~$y$ lan nyukupi $R(\vec x, z)$,
yen ana, utawa ngasilake $y$ yen ora ana, uga rekursif
primitif. Kita bakal nulis fungsi~$m_R$ minangka
\[
\bmin{z < y}{R(\vec{x}, z)},
\]
\end{prop}

\begin{proof}
Ora ana~$z < 0$ kang nyukupi $R(\vec{x}, z)$ amarga
pancen ora ana $z < 0$. Mula $m_R(\vec x, 0) = 0$.

Yen watese awujud $y + 1$, ana telung kasus:
\begin{enumerate}
\item Ana $z < y$ kang nyukupi $R(\vec{x}, z)$;
mula $m_R(\vec{x}, y+1) = m_R(\vec{x}, y)$.
\item Ora ana~$z<y$ mangkono, nanging $R(\vec{x}, y)$
bener; mula $m_R(\vec{x}, y+1) = y$.
% OLPL-124: Sumber nggunakake vektor z kanggo input m_R ing kasus
% katelu; proposisi, rong kasus sadurunge lan rekurensi nganggo vektor x.
\item Ora ana $z < y+1$ kang nyukupi $R(\vec{x}, z)$;
mula $m_R(\vec{x}, y+1) = y+1$.
\end{enumerate}
Mula kita bisa netepake $m_R(\vec x, y)$ kanthi rekursi
primitif kaya mangkene:
\begin{align*}
m_R(\vec{x}, 0) & = 0\\
m_R(\vec{x}, y+1) & =
\begin{cases}
m_R(\vec{x}, y) & \text{yen $m_R(\vec x, y) \neq y$}\\
y & \text{yen $m_R(\vec x, y) = y$ lan $R(\vec{x}, y)$}\\
y+1 & \text{yen ora.}
\end{cases}
\end{align*}
Cathetan koreksi sumber OLPL-124: parameter ing kasus katelu yaiku vektor x, dudu vektor z. Cathetan koreksi sumber SOL6-F037: ukara pambuka rumus iki netepake fungsi kanggo kabeh wates, dudu mung nilai wiwitan ing nol.
Gatekna yen ana $z<y$ kang nyukupi $R(\vec x, z)$
yen lan mung yen $m_R(\vec x,
y) \neq y$.
\end{proof}

\begin{prob}
Upamana $R(\vec x, z)$ rekursif primitif. Netepna fungsi
$m'_R(\vec{x}, y)$ kang ngasilake~$z$ paling cilik kang
kurang saka~$y$ lan nyukupi $R(\vec{x}, z)$ yen ana,
utawa ngasilake $0$ yen ora ana, kanthi rekursi
primitif saka~$\Char{R}$.
\end{prob}

\end{document}
